% from ega3-1-fr.tex
\oldpage[III]{158}

Soit $\mathcal L$ un $\mathcal O_{\mathfrak X}$-Module inversible, et supposons
que $\mathcal L_0=\mathcal L/\mathfrak J\mathcal L$ soit un
$\mathcal O_{X_0}$-Module ample. Alors $\mathfrak X$ est algébrisable, et si
$X$ est un $S$-schéma propre tel que $\mathfrak X$ soit isomorphe à
$\widehat X$, il existe un $\mathcal O_X$-Module ample $\mathcal M$ tel que
$\mathcal L$ soit isomorphe à $\widehat{\mathcal M}$ (ce qui entraîne que
$X$ est projectif sur $S$).
\end{theorem}

\begin{proof}
Appliquons (\hyperref[III.5.2.3-fr]{5.2.3}) à
$\mathcal F=\mathcal O_{\mathfrak X}$~: il existe donc un entier $n_0$ tel que
pour $n\geq n_0$, l'homomorphisme canonique
$\Gamma(\mathfrak X,\mathcal L^{\otimes n})\to
\Gamma(X_0,\mathcal L_0^{\otimes n})$ soit surjectif. On peut supposer
$n\geq n_0$ choisi assez grand pour que $\mathcal L_0^{\otimes n}$ soit très
ample pour $S_0$ (II, 4.5.10). Comme le morphisme $f_0:X_0\to S_0$ est propre,
$\Gamma(X_0,\mathcal L_0^{\otimes n})$ est un $A$-module de type fini (3.2.1),
donc il existe un sous-$A$-module de type fini $E$ de
$\Gamma(\mathfrak X,\mathcal L^{\otimes n})$ dont l'image dans
$\Gamma(X_0,\mathcal L_0^{\otimes n})$ soit ce dernier module tout entier.
Cela étant, pour tout $k\geq0$, considérons l'homomorphisme
\[
u_k:E/\mathfrak J^{k+1}E\longrightarrow
\Gamma(X_k,\mathcal L_k^{\otimes n})
\]
déduit de l'injection canonique
$E\to\Gamma(\mathfrak X,\mathcal L^{\otimes n})$. Notons que
$(f_k)_*(\mathcal L_k^{\otimes n})$ est quasi-cohérent, et comme
$\Gamma(S_k,(f_k)_*(\mathcal L_k^{\otimes n}))
=\Gamma(X_k,\mathcal L_k^{\otimes n})$, $u_k$ définit un homomorphisme
\[
\widetilde u_k:(E/\mathfrak J^{k+1}E)^\sim
\longrightarrow(f_k)_*(\mathcal L_k^{\otimes n}),
\]
et par suite aussi un homomorphisme
\[
\widetilde u_k^\sharp:
f_k^*((E/\mathfrak J^{k+1}E)^\sim)\longrightarrow
\mathcal L_k^{\otimes n}.
\]
D'ailleurs, si on pose
$\mathcal G_k=f_k^*((E/\mathfrak J^{k+1}E)^\sim)$, on a
$\mathcal G_0=\mathcal G_k/\mathfrak J\mathcal G_k$ (I, 9.1.5), donc
$\widetilde u_0^\sharp:\mathcal G_k/\mathfrak J\mathcal G_k
\to\mathcal L_k^{\otimes n}/\mathfrak J\mathcal L_k^{\otimes n}$ se déduit
de $\widetilde u_k^\sharp$ par passage aux quotients. Or, par définition de
$E$, $\widetilde u_0^\sharp$ n'est autre que l'homomorphisme canonique
\[
\sigma:f_0^*((f_0)_*(\mathcal L_0^{\otimes n}))\longrightarrow
\mathcal L_0^{\otimes n},
\]
et l'hypothèse que $\mathcal L_0^{\otimes n}$ est très ample entraîne que
$\widetilde u_0^\sharp$ est surjectif (II, 4.4.3)~; on déduit alors du lemme
de Nakayama que chacun des $\widetilde u_k^\sharp$ est aussi surjectif.
Chacun des $\widetilde u_k^\sharp$ définit donc (II, 4.2.2) un
$S_k$-morphisme
\[
g_k:X_k\longrightarrow P_k=\mathbf P(E/\mathfrak J^{k+1}E),
\]
et comme $P_h=P_k\times_{S_k}S_h$ pour $h\leq k$ en vertu de (II, 4.1.3),
$(g_k)$ est un $(S_k)$-système inductif adique (I, 10.12.2) en vertu des
relations entre les $\widetilde u_k^\sharp$ et de (II, 4.2.10). Les $g_k$
définissent donc un $\mathfrak S$-morphisme de schémas formels
$g:\mathfrak X\to\mathfrak P$, où $\mathfrak P$ est la limite inductive du
système $(P_k)$, ou encore le complété $\widehat P$, où $P=\mathbf P(E)$. De
plus, l'hypothèse que $\mathcal L_0^{\otimes n}$ est très ample entraîne que
$g_0$ est une immersion fermée (II, 4.4.3)~; on en conclut que $g$ est une
immersion fermée de schémas formels (4.8.10), donc $\mathfrak X$ est
algébrisable (5.1.8). Le fait que $\mathcal L$ soit isomorphe au complété
$\widehat{\mathcal M}$ d'un $\mathcal O_X$-Module inversible résulte alors du
th. d'existence (5.1.6). En outre, $\mathcal L^{\otimes n}$ est alors le
complété de $\mathcal M^{\otimes n}$ (I, 10.8.10), et les homomorphismes
$\widetilde u_k^\sharp$ définissent un homomorphisme bien déterminé
\[
v:f^*(\widetilde E)\longrightarrow\mathcal M^{\otimes n}
\]
(5.1.7)~; d'ailleurs, comme $\widetilde u_k^\sharp$ est surjectif, il en est de
même de $\widehat v$ (I, 10.11.5), donc de $v$ (5.1.3)~; en outre, le morphisme
$r:X\to P$ défini par $v$ (II, 4.2.2) a pour prolongement aux complétés $g$,
et comme $g$ est une immersion fermée, il en est de même de $r$, par (5.1.8)
et (5.4.1)~; on en conclut que $\mathcal M^{\otimes n}$ est très ample
(II, 4.4.6) et $\mathcal M$ est ample (II, 4.5.10).
\end{proof}

\begin{remark}[5.4.6]
\label{III.5.4.6-fr}
Soient $A$ un anneau adique noethérien,
$\mathfrak S=\operatorname{Spf}(A)$,
$f:\mathfrak X\to\mathfrak S$ un morphisme propre de schémas formels. Soit
$\mathcal N$ l'Idéal cohérent de $\mathcal O_{\mathfrak X}$ tel que pour tout
ouvert formel affine $U$ de $\mathfrak X$, $\Gamma(U,\mathcal N)$ soit le
nilradical de $\Gamma(U,\mathcal O_{\mathfrak X})$~; l'existence de cet Idéal
résulte facilement de (I, 10.10.2) et du fait que tout homomorphisme d'anneaux
$B\to C$ applique le nilradical de $B$ dans celui de $C$. Soit
$\mathfrak X'$ le sous-schéma formel fermé de $\mathfrak X$ défini par
$\mathcal N$ (I, 10.14.2)~; il serait intéressant de savoir si, lorsque
$\mathfrak X'$ est algébrisable, $\mathfrak X$ lui-même est algébrisable. On
parviendrait sans doute à une solution de ce problème si on savait classifier
(par exemple au moyen d'invariants de nature

